% !TeX root = ../colloquium-2.tex

\hiddensection{Определения}

%% def 1
% Непрерывные на множестве функции. Теорема о нуле непрерывной на отрезке функции.
% Определение ограниченной на множестве функции.
% % !TeX root = ../colloquium-2.tex

\phantomsection\label{sec_Definition01}
\begin{definition}[Непрерывные на множестве функции. 
      Теорема о нуле непрерывной на отрезке функции.
      Определение ограниченной на множестве функции.]
    \subsection*{Непрерывные на множестве функции.}
    \begin{note}[purple]
        Функция $f$, определённая на множестве $E$, называется непрерывной на $E$,
        если она непрерывна в каждой точке $E$. 
    \end{note}
    \subsection*{Теорема о нуле непрерывной на отрезке функции.}
    \begin{note}[purple]
        Пусть функция $f \in C([a, b])$ и $f(a) \cdot f(b) < 0$. Тогда 
        $\exists c \in (a, b) \; : \; f(c) = 0$. 
    \end{note}
    \subsection*{Определение ограниченной на множестве функции.}
    \begin{note}[purple]
        
    \end{note}
\end{definition}
% !TeX root = ../colloquium-2.tex

\phantomsection\label{sec_Definition01}
\begin{definition}[Непрерывные на множестве функции. 
      Теорема о нуле непрерывной на отрезке функции.
      Определение ограниченной на множестве функции.]
    \subsection*{Непрерывные на множестве функции.}
    \begin{note}[purple]
        Функция $f$, определённая на множестве $E$, называется непрерывной на $E$,
        если она непрерывна в каждой точке $E$. 
    \end{note}
    \subsection*{Теорема о нуле непрерывной на отрезке функции.}
    \begin{note}[purple]
        Пусть функция $f \in C([a, b])$ и $f(a) \cdot f(b) < 0$. Тогда 
        $\exists c \in (a, b) \; : \; f(c) = 0$. 
    \end{note}
    \subsection*{Определение ограниченной на множестве функции.}
    \begin{note}[purple]
        
    \end{note}
\end{definition}

%% def 2
% 1-я теорема Вейерштрасса.
% 2-я теорема Вейерштрасса.
% Теорема Коши о непрерывной функции.
% !TeX root = ../colloquium-2.tex

% 1-я теорема Вейерштрасса.
% 2-я теорема Вейерштрасса.
% Теорема Коши о непрерывной функции.

\phantomsection\label{sec_Definition02}
\begin{definition}[
    1-я теорема Вейерштрасса.
    2-я теорема Вейерштрасса.
    Теорема Коши о непрерывной функции.
]
    \subsection*{1-я теорема Вейерштрасса.}
    \begin{note}[purple]
        Пусть функция $f$ определена и непрерывна на отрезке $[a, b]$. 
        Тогда она ограничена на этом отрезке.
    \end{note}

    \subsection*{2-я теорема Вейерштрасса.}
    \begin{note}[purple]
        Пусть функция $f$ определена и непрерывна на отрезке $[a, b]$. 
        Тогда существуют такие точки $x_1, \; x_2 \in [a, b]$, что 
        $$
            f(x_1) = m = \inf_{a \leq x \leq b} f(x), \ f(x_2) = M = \sup_{a \leq x \leq b} f(x)
        $$
    \end{note}

    \subsection*{Теорема Больцано-Коши о непрерывной функции.}
    \begin{note}[purple]
        Пусть функция $f$ определена и непрерывна на отрезке 
        $[a, b]$, $m = \displaystyle{\inf_{a \leq x \leq b} f(x)}$,
        $M = \displaystyle{\sup_{a \leq x \leq b} f(x)}$. \\

        Тогда для любого числа $C \in [m, \; M]$ существует такая точка 
        $c \in [a, \; b]$, что $f(c) = C$. 
    \end{note}
\end{definition}

%% def 3
% Определение равномерно непрерывной функции. 
% Теорема Гейне-Кантора о равномерной непрерывности.
% !TeX root = ../colloquium-2.tex

% Определение равномерно непрерывной функции. 
% Теорема Гейне-Кантора о равномерной непрерывности.

\phantomsection\label{sec_Definition03}
\begin{definition}[Определение равномерно непрерывной функции. 
                   Теорема Гейне-Кантора о равномерной непрерывности.]
    \subsection*{Определение равномерно непрерывной функции.}
    \begin{note}[purple]
        Функция $f$, определённая на множестве $E$, называется
        равномерно непрерывной на этом множестве, если для любого числа
        $\epsilon > 0$ существует такое $\Delta > 0$, что при всех
        таких $x_1, \, x_2 \in E$, что $|x_1 - x_2| < \Delta$, 
        выполнено неравенство
        $$
            |f(x_1) - f(x_2)| < \epsilon
        $$
    \end{note}
    \subsection*{Теорема Гейне-Кантора о равномерной непрерывности.}
    \begin{note}[purple]
        Функция $f$, непрерывная на отрезке $[a, \; b]$, 
        равномерно непрерывна на этом отрезке.
    \end{note}
\end{definition}

%% def 4
% Определение обратной функции.
% Критерий непрерывности монотонной функции.
% Теорема об обратной функции.
% !TeX root = ../colloquium-2.tex

% Определение обратной функции.
% Критерий непрерывности монотонной функции.
% Теорема об обратной функции.

\phantomsection\label{sec_Definition04}
\begin{definition}[Определение обратной функции. 
                   Критерий непрерывности монотонной функции.
                   Теорема об обратной функции.]
    \subsection*{Определение обратной функции.}
    \begin{note}[purple]
        Пусть функция $f \; : \; E \rightarrow D$ осуществляет биекцию между $E$ и $D$.
        Если каждому $y \in D$ поставить в соответствие то $x \in E$, для которого $f(x) = y$, 
        то тем самым будет определена функция, отображающая множество $D$ во множество $E$. \\

        Эта функция будет называться \textbf{обратной} для функции $f$ и обозначается $f^{-1}$. \\
        Таким образом, $f^{-1} \; : \; D \rightarrow E$. 
    \end{note}

    \subsection*{Критерий непрерывности монотонной функции.}
    \begin{note}[purple]
        Монотонная на отрезке $[a, \, b]$ функция $f$ непрерывна на этом отрезке тогда и только тогда, 
        когда множеством её значений является отрезок с концами $f(a)$ и $f(b)$. 
    \end{note}

    \subsection*{Теорема об обратной функции.}
    \begin{note}[purple]
        Пусть функция $f$ непрерывна и строго монотонна (то есть возрастает или убывает) на 
        отрезке $[a, \, b]$. \\
        Тогда функция $f$ имеет обратную функцию $f^{-1}$, определённую на отрезке с концами 
        $f(a)$ и $f(b)$, причём $f^{-1}$ строго монотонна и непрерывна на отрезке с концами
        $f(a)$ и $f(b)$ и характер монотонности функций $f$ и $f^{-1}$ одинаковый. 
    \end{note}
\end{definition}


%% def 5
% Определение дифференцируемой функции. 
% Определение дифференциала.
% Линейная функция.
% !TeX root = ../colloquium-2.tex

% Определение дифференцируемой функции. 
% Определение дифференциала.
% Линейная функция.

\phantomsection\label{sec_Definition05}
\begin{definition}[Определение дифференцируемой функции. 
                   Определение дифференциала. 
                   Линейная функция.]
    \subsection*{Определение дифференцируемой функции.}
    \begin{note}[purple]
        Функция $f$, определённая в некоторой окрестности точки $a$,
        называется \textbf{дифференцируемой} в точке $a$, если 
        существуют такие число $A$ и функция $\alpha$, что 
        для всех $h$ из некоторой проколотой окрестности нуля 
        выполнено равенство 
        $$f(a + h) - f(a) = Ah + \alpha(h)h,$$ 
        где $\displaystyle{\lim_{h \rightarrow 0} \alpha(h) = 0}$. \\
        При этом $A$ и $\alpha$ зависят и от точки $a$, поэтому часто
        равенство записывают в виде 
        $$ f(a + h) - f(a) = A(a)h + \alpha(a, h)h $$
    \end{note}

    \subsection*{Определение дифференциала.}
    \begin{note}[purple]
        Функция $h \mapsto Ah$ называется \textbf{дифференциалом} 
        функции $f$ в точке $a$. Она обозначается $df(a)$ или
        $df | _{x = a}$, то есть $df(a)(h) = df(h)|_{x = a} = Ah$.
    \end{note}

    \subsection*{Линейная функция.}
    \begin{note}[purple]
        Линейной функцией называется функция вида $f = kx + b$, 
        где $x$ - независимая переменная, $k$ и $b$ - числа. 
    \end{note}
\end{definition}


%% def 6
% Определение производной. 
% Связь дифференцируемости и производной.
% Непрерывность дифференцируемой функции.
% !TeX root = ../colloquium-2.tex

% Определение производной. 
% Связь дифференцируемости и производной.
% Непрерывность дифференцируемой функции.

\phantomsection\label{sec_Definition06}
\begin{definition}[Определение производной.
                   Связь дифференцируемости и производной.
                   Непрерывность дифференцируемой функции.]
    \subsection*{Определение производной.}
    \begin{note}[purple]
        Если существует предел 
        $\displaystyle{\lim_{h \rightarrow 0}} \frac{f(a + h) - f(a)}{h}$, 
        то он называется производной функции $f$ в точке $a$. \\
        Другая форма записи предела:
        $$ \lim_{x \rightarrow a} \frac{f(x) - f(a)}{x - a} $$
        Производная функции $f$ в точке $a$ обозначается символом 
        $f'(a)$ или $\frac{df}{dx}(a)$. 
    \end{note}

    \subsection*{Связь дифференцируемости и производной.}
    \begin{note}[purple]
        Функция $f$ дифференцируема в точке $a$ тогда и только тогда, когда 
        в точке $a$ существует производная этой функции $f'(a)$. При этом 
        $df(h) | _{x = a} = f'(a)h$.
    \end{note}

    \subsection*{Непрерывность дифференцируемой функции.}
    \begin{note}[purple]
        Пусть функция $f$ дифференцируема в точке $a$. Тогда $f$ непрерывна в точке $a$. 
    \end{note}
\end{definition}


%% def 7
% Формулировка правил дифференцирования.
% !TeX root = ../colloquium-2.tex

% Формулировка правил дифференцирования.

\phantomsection\label{sec_Definition07}
\begin{definition}[Формулировка правил дифференцирования.]
    Пусть функции $f$ и $g$ дифференцируемы в точке $a$. \\
    Тогда:
    \begin{enumerate}
        \item \textbf{Свойство линейности:} \\
        Функция $\alpha f + \beta g$ дифференцируема в точке $a$ при любых 
        $\alpha, \; \beta \in \mathbb{R}$ и 
        $$ (\alpha f + \beta g)' (a) = \alpha f'(a) + \beta g' (a) $$

        \item \textbf{Правило Лейбница:} \\
        $f \cdot g$ дифференцируема в точке $a$, а также $(f \cdot g)'(a) = f'(a)g(a) + f(a)g'(a)$

        \item Если $g \neq 0$ в некоторой окрестности точки $a$, то $\frac{f}{g}$ дифференцируема
        в точках $a$ и $\big( \frac{f}{g} \big)' (a) = \frac{f'(a)g(a) - g'(a)f(a)}{g^2(a)}$.
    \end{enumerate}
\end{definition}


%% def 8
% Инвариантность формы первого дифференциала. 
% Теорема о производной сложной функции.
% !TeX root = ../colloquium-2.tex

% Инвариантность формы первого дифференциала. 
% Теорема о производной сложной функции.

\phantomsection\label{sec_Definition08}
\begin{definition}[Инвариантность формы первого дифференциала.
                   Теорема о производной сложной функции.]
    \subsection*{Инвариантность формы первого дифференциала.}
    \begin{note}[purple]
        Вне зависимости от того, является ли $y$ независимой и переменной
        функцией, форма (то есть вид) первого дифференциала внешне не меняется.
    \end{note}

    \subsection*{Теорема о производной сложной функции.}
    \begin{note}[purple]
        Пусть функция $f$ дифференцируема в точке $a$, а функция $g$ дифференцируема в точке
        $f(a)$. Тогда функция $g \circ f$ дифференцируема в точке $a$ и 
        $(g \circ f)'(a) = g'(f(a)) \cdot f'(a)$. 
    \end{note}
\end{definition}


%% def 9
% Теорема о производной обратной фукнции.
% !TeX root = ../colloquium-2.tex

% Теорема о производной обратной фукнции.

\phantomsection\label{sec_Definition09}
\begin{definition}[Теорема о производной обратной функции.]
    Пусть $f$ - непрерывная и строго монотонная функция, отображающая интервал $I$ 
    в интервал $J$. \\
    Пусть также $f$ дифференцируема в точке $a \in I$ и $f'(a) \neq 0$. \\
    Тогда обратная функция $f^{-1}$ дифференцируема в точке $b = f(a)$ и 
    $(f^{-1})'(b) = \frac{1}{f'(a)}$.
\end{definition}


%% def 10
% Таблица производных.
% !TeX root = ../colloquium-2.tex

% Таблица производных.

\phantomsection\label{sec_Definition10}
\begin{definition}[Таблица производных.]
    \begin{enumerate}
        % 1
        \item $$ C' = 0 \; \forall C = const $$
        % 2
        \item $$ (x^{\alpha})' = \alpha x^{\alpha - 1} \; (x > 0) $$
        % 3
        \item $$ (e^x)' = e^x $$
        % 4
        \item $$ (a^x)' = a^x \ln a(a > 0, \; a \neq 1) $$
        % 5
        \item $$ (\ln x)' = \frac{1}{x} \; (x > 0) $$
        % 6
        \item $$ (\log_a x)' = \frac{1}{x \ln{a}} \; (x > 0, \; a > 0, \; a \neq 1) $$
        % 7
        \item $$ (\sin x)' = \cos x $$
        % 8
        \item $$ (\cos x) = - \sin x $$
        % 9
        \item $$ (\tg x)' = \frac{1}{\cos^2 x} = tg^2 x + 1 \; (x \neq \frac{\pi}{2} + \pi k, \; k \in \mathbb{Z}) $$
        % 10
        \item $$ (\ctg x)' = \frac{1}{\sin^2 x} = - ctg^2 x $$
        % 11
        \item $$ (\arcsin x)' = \frac{1}{\sqrt{1 - x^2}} \; (-1 < x < 1) $$
        % 12
        \item $$ (\arccos x)' = - \frac{1}{\sqrt{1 - x^2}} \; (-1 < x < 1) $$
        % 13
        \item $$ (\arctg x)' = \frac{1}{1 + x^2} $$
        % 14
        \item $$ (\arcctg x)' = - \frac{1}{1 + x^2} $$
        % 15
        \item $$ (sh x)' = ch x $$
        % 16
        \item $$ (ch x)' = sh x $$
    \end{enumerate}
\end{definition}


%% def 11
% Определение локального экстремума и формулировка теоремы Ферма.
% !TeX root = ../colloquium-2.tex

% Определение локального экстремума и формулировка теоремы Ферма.

\phantomsection\label{sec_Definition11}
\begin{definition}[Определение локального экстремума и формулировка теоремы Ферма.]
    \subsection*{Определение локального экстремума.}
    \begin{note}[purple]
        Пусть $\delta > 0$, $U_{\delta}(a)$ - $\delta$-окрестность точки $a$. \\
        Если функция $f$ определена на множестве $E$, точка $a \in E$ и при всех 
        $x \in U_{\delta}(a) \cap E$, выполнено неравенство $f(x) \leq f(a)$, 
        то $a$ называется точкой локального минимума функции $f$. \\
        Если неравенство строгое, то $a$ называется точкой строгого локального максимума
        или строго локального минимума соответственно. Точки \textbf{максимума} и \textbf{минимума} называются
        точками \textbf{локального экстремума} (если неравенства нестрогие), и \textbf{точками строгого 
        локального экстремума}.
    \end{note}

    \subsection*{Теорема Ферма.}
    \begin{note}[purple]
        Пусть функция $f$ определена на интервале $(a, b)$, дифференцируема в точке $c \in (a, b)$
        и имеет в точке $c$ локальный экстремум. \\
        Тогда $f'(c) = 0$. 
    \end{note}
\end{definition}


%% def 12
% Формулировки теорем Ролля и Лагранжа.
% !TeX root = ../colloquium-2.tex

% Формулировки теорем Ролля и Лагранжа.

\phantomsection\label{sec_Definition12}
\begin{definition}[Формулировки теорем Ролля и Лагранжа.]
    \subsection*{Теорема Ролля.}
    \begin{note}[purple]
        Пусть:
        \begin{enumerate}
            \item Функция $f$ определена и непрерывна на отрезке $[a,\; b]$
            \item Функция $f$ дифференцируема на интервале $(a, \; b)$
            \item $f(a) = f(b)$
        \end{enumerate}
        Тогда существует такая точка $c \in (a, \; b)$, что $f'(c) = 0$. 
    \end{note}

    \subsection*{Теорема Лагранжа.}
    \begin{note}[purple]
        Пусть:
        \begin{enumerate}
            \item Функция $f$ определена и непрерывна на отрезке $[a,\; b]$
            \item Функция $f$ дифференцируема на интервале $(a, \; b)$
        \end{enumerate}
        Тогда существует такая точка $c \in (a, \; b)$, что $f(b) - f(a) = f'(c)(b - a)$.
    \end{note}
\end{definition}


%% def 13
% Формулировки следствий теоремы Лагранжа.
% Формулировка теоремы Коши.
% !TeX root = ../colloquium-2.tex

% Формулировки следствий теоремы Лагранжа.
% Формулировка теоремы Коши.

\phantomsection\label{sec_Definition13}
\begin{definition}[Формулировки следствий теоремы Лагранжа.
                   Формулировка теоремы Коши.]
    \subsection*{Формулировки следствий теоремы Лагранжа.}
    \begin{note}[amber]
        \subsubsection*{Следствие 1.}
        \begin{note}[purple]
            Пусть определена и дифференцируема на интервале $(a, \; b)$
            и $f'(x) = 0 \; \forall x \in (a, \; b)$. Тогда функция $f$ постоянна на $(a, \; b)$. 
        \end{note}

        \subsubsection*{Следствие 2.}
        \begin{note}[purple]
            Пусть функция $f$ определена и дифференцируема на интервале $(a, \; b)$
            \begin{enumerate}
                \item Функция $f$ не убывает на этом интервале тогда и только тогда, 
                когда $$ f'(x) \geq 0 \; \forall x \in (a, \; b) $$
                \item Функция $f$ не возрастает на этом интервале тогда и только тогда, 
                когда $$ f'(x) \leq 0 \; \forall x \in (a, \; b) $$
            \end{enumerate}
        \end{note}
    \end{note}

    \subsection*{Формулировка теоремы Коши.}
    \begin{note}[purple]
        Пусть функции $f$ и $g$:
        \begin{enumerate}
            \item определены и непрерывны на отрезке $[a,\; b]$
            \item дифференцируемы на интервале $(a, \; b)$
            \item $g'(x) \neq 0 \; \forall x \in (a, \; b)$
        \end{enumerate}
        Тогда существует такая точка $c \in (a, \; b)$, что 
        $$ \frac{f(b) - f(a)}{g(b) - g(a)} = \frac{f'(c)}{g'(c)} $$
    \end{note}
\end{definition}


%% def 14
% Формулировки первого и второго правил Лопиталя.
% !TeX root = ../colloquium-2.tex

% Формулировки первого и второго правил Лопиталя.

\phantomsection\label{sec_Definition14}
\begin{definition}[Формулировки первого и второго правил Лопиталя.]
    \subsection*{Формулировка первого правила Лопиталя, 
                 неопределённость вида $\frac{0}{0}$.}
    \begin{note}[purple]
        Пусть:
        \begin{enumerate}
            \item Функции $f$ и $g$ определены и дифференцируемы на интервале $(a, \; b)$
            \item $ \ds \lim_{x \rightarrow b-} f(x) = \lim_{x \rightarrow b-} g(x) = 0 $
            \item $g'(\xi) \neq 0 \; \forall \xi \in (a, \; b)$
            \item Существует конечный или бесконечный предел 
            $\ds \lim_{x \rightarrow b-} \frac{f'(x)}{g'(x)}$
        \end{enumerate}
        Тогда справедливо равенство 
        $\ds \lim_{x \rightarrow b-} \frac{f(x) }{g(x) } = 
                       \lim_{x \rightarrow b-} \frac{f'(x)}{g'(x)} $
    \end{note}

    \subsection*{Формулировка второго правила Лопиталя,
                 неопределённость вида $\frac{\infty}{\infty}$.}
    \begin{note}[purple]
        \begin{enumerate}
            \item Функции $f$ и $g$ определены и дифференцируемы на интервале $(a, \; b)$
            \item $\lim_{x \rightarrow b-} g(x) = \infty$
            \item $g'(\xi) \neq 0 \; \forall \xi \in (a, \; b)$
            \item Существует конечный или бесконечный предел 
            $\ds \lim_{x \rightarrow b-} \frac{f'(x)}{g'(x)}$
        \end{enumerate}
        Тогда справедливо равенство $\ds \lim_{x \rightarrow b-} \frac{f(x)} {g(x)} = 
                                         \lim_{x \rightarrow b-} \frac{f'(x)}{g'(x)}$
    \end{note}
\end{definition}


%% def 15
% Определение производной порядка больше или равно 2. 
% Формулировка обобщённого правила Лейбница.
% !TeX root = ../colloquium-2.tex

% Определение производной порядка больше или равно 2. 
% Формулировка обобщённого правила Лейбница.

\phantomsection\label{sec_Definition15}
\begin{definition}[
    Определение производной порядка больше или равно 2.
    Формулировка обобщённого правила Лейбница.
]
    \subsection*{Определение производной порядка больше или равно 2.}
    Пусть функция $f$ дифференцируема в каждой точке множества $E$. \\
    Тогда её производная $f'$ определена на $E$ и можно ставить вопрос о её 
    дифференцируемости. \\
    Отметим, что производная функции $f$ нулевого порядка называется сама функция $f$.
    Производная порядка $n$ определяется индуктивно.
    \begin{note}[purple]
        Если у функции $f$ в некоторой окрестности точки $a$ существует производная
        $n - 1$-го порядка $f^{(n - 1)}$ (то есть эта производная представляет собой
        определённую в некоторой окрестности точки $a$ функцию) и эта производная 
        дифференцируема в точке $a$, то производная функции $f$ порядка $n$ в точке $a$ - 
        это производная функции $f^{(n - 1)}$ в точке $a \; : \; f^{(n)}(a) := (f^{(n - 1)})'(a)$. \\

        \textbf{В терминах пределов:}
        $$ 
            f^{(n)}(a) = \lim_{h \rightarrow 0} \frac{f^{(n - 1)}(a + h) - f^{(n - 1)}(a)}{h}
        $$

        Ещё раз отметим, что производная порядка $n - 1$ определена в окрестности точки $a$, 
        иначе символ $f^{(n - 1)} (a + h)$ не имеет смысла.
    \end{note}

    \subsection*{Формулировка обобщённого правила Лейбница.}
    \begin{note}[purple]
        Пусть функции $u$ и $v$ $n$ раз дифференцируемы в точке $a$. 
        Тогда их произведение также $n$ раз дифференцируемо в точке $a$ и
        $$ (u \cdot v)^{(n)} (a) = \sum_{k = 0}^{n} C_n^k u^{(k)} (a) v^{(n - k)}(a) $$
    \end{note}
\end{definition}


%% def 16
% Определение многочлена Тейлора и формулировка теоремы о разложении с
%       остаточным членом в форме Пеано.
% !TeX root = ../colloquium-2.tex

% Определение многочлена Тейлора и формулировка теоремы о 
% разложении с остаточным членом в форме Пеано.

\phantomsection\label{sec_Definition16}
\begin{definition}[
    Определение многочлена Тейлора и формулировка теоремы о разложении
    с остаточным членом в фореме Пеано.
]
    Если разницу между функцией и её многочленом Тейлора записать в виде формы Пеано,
    то никакой дополнительной информации, кроме той, что это разница есть
    бесконечно малая величина при $x \rightarrow a$, мы не имеем. Иногда, например, 
    при оценке погрешости важно иметь более подробную информацию о поведении выражения
    $$ R_n (x; \; f, \; a) := f(x) - T_n (x; \; f, \; a) $$
\end{definition}


%% def 17
% Формулировка о разложении по формуле Тейлора с остаточным членом в общей форме.
% !TeX root = ../colloquium-2.tex

% Формулировка о разложении по формуле Тейлора с остаточным членом в общей форме.

\phantomsection\label{sec_Definition17}
\begin{definition}[
    Формулировка о разложении по формуле Тейлора с остаточным членом в общей форме.
]
    Пусть в каждой точке отрезка с концами $a$ и $x$ функция $f$ имеет $n$ непрерывных
    производных, а в каждой точке интервала с концами $a$ и $x$ $f$ дифференцируема $n + 1$ раз.
    Пусть также функция $g$ непрерывна на отрезке с концами $a$ и $x$ и дифференцируема в каждой 
    точке интервала с концами $a$ и $x$, причём $g'(\xi) \neq 0$ в любой точке $\xi$ этого 
    интервала. Тогда 
    $$ R_n (x; \; f, \; a) = \frac{f^{(n + 1)} (c)}{n! g'(c)} (g(x) - g(a))(x - c)^n, $$
    где $c$ - какая-то точка из интервала с концами $a$ и $x$.
\end{definition}

%% def 18
% Формулировка достаточного условия экстремума.
% Определение Определение выпуклой функции.
% Формулировка критерия выпуклости.
% !TeX root = ../colloquium-2.tex

% Формулировка достаточного условия экстремума.
% Определение Определение выпуклой функции.
% Формулировка критерия выпуклости.

\phantomsection\label{sec_Definition18}
\begin{definition}[
    Формулировка достаточного условия экстремума.
    Определение выпуклой функции.
    Формулировка критерия выпуклости.
]
    \subsection*{Формулировка достаточного условия экстремума.}
    \begin{note}[purple]
        Пусть функция $f$ определена в некоторой окрестности точки $a$ и дифференцируема $n$ 
        раз в точке $a$, причём 
        $$ f'(a) = ... = f^{(n - 1)}(a) = 0, $$
        а $f^{(n)}(a) \neq 0$. Тогда если $n = 2k$ и $f^{(n)}(a) < 0$, то $a$ - точка максимума,
        если $n = 2k$ и $f^{(n)}(a) > 0$, то $a$ - точка минимума, а если $n = 2k - 1$, то
        экстремума нет.
    \end{note}

    \subsection*{Определение выпуклой функции.}
    \begin{note}[purple]
        \subsubsection*{Определение выпуклости}
        Функция $f$, определённая на интервале $(a, \; b)$, называется выпуклой 
        (\textit{выпуклой вниз}) на этом интервале, если для любых точек $x_1, x_2 \in (a, \; b)$
        и любого числа $\lambda \in [0, \; 1]$ выполнено неравенство
        $$ f(\lambda x_1 + (1 - \lambda)x_2) \leq \lambda f(x_1) + (1 - \lambda) f(x_2) $$

        Если при тех же условиях выполнено противоположное неравенство, то функцию $f$ называют
        вогнутой (\textit{выпуклой вверх}) на интервале $(a, \; b)$/

        \subsubsection*{Геометрическое определение выпуклости}
        Всякая хорда, соединяющая две точки граифика функции, лежит не ниже графика функции,
        а определение вогнутости означает, что любая хорда лежит не выше графика функции.
    \end{note}

    \subsection*{Формулировка критерия выпуклости.}
    \begin{note}[purple]
        Функция $f$ выпукла на интервале $(a, \; b)$ тогда и только тогда, когда для любых
        таких чисел $x_1, \; x_2, \; y \in (a, \; b)$, что $x_1 < y < x_2$, выполнено 
        неравенство 
        $$ 
        \frac{f(y) - f(x_1)}{y - x_1} \leq 
        \frac{f(x_2) - f(y)}{x_2 - y}
        $$
        Геометрический смысл заключается в том, что тангенсы углов наклонов секущих не убывают
        с увеличением абсцисс точек пересечения с графиком функции $f$. 
    \end{note}
\end{definition}

%% def 19
% Формулировки критериев выпуклости в терминах первой и второй производных.
% Формулировка теоремы о взаимном расположении выпуклой функции и касательной. 
% Неравенство Йенсена.
% !TeX root = ../colloquium-2.tex

% Формулировки критериев выпуклости в терминах первой и второй производных.
% Формулировка теоремы о взаимном расположении выпуклой функции и касательной. 
% Неравенство Йенсена.

\phantomsection\label{sec_Definition19}
\begin{definition}[
    Формулировки критериев выпуклости в терминах первой и второй производных.
    Формулировка теоремы о взаимном расположении выпуклой функции и касательной. 
    Неравенство Йенсена.
]
    \subsection*{Формулировки критериев выпуклости в терминах первой и второй производных.}
    \begin{note}[purple]
        Пусть функция $f$ дифференцируема на интервале $(a, \; b)$. Тогда выпуклость $f$ 
        на этом интервале равносильна тому, что функция $f'$ не убывает на $(a, \; b)$.
    \end{note}

    \subsection*{Формулировка теоремы о взаимном расположении выпуклой функции и касательной.}
    \begin{note}[purple]
        Пусть функция $f$ дважды дифференцируема на интервале $(a, \; b)$. Тогда выпуклость $f$
        на этом интервале равносильна тому, что функция $f''$  неотрицательна на $(a, \; b)$. 
    \end{note}

    \subsection*{Неравенство Йенсена.}
    \begin{note}[purple]
        Пусть функция $f$ выпукла на интервале $(a, \; b)$. Тогда при любых 
        $x_1, \; x_2, \; ..., \; x_n \in (a, \; b)$ и всех таких 
        $\lambda_1 \geq 0, \; \lambda_2 \geq 0, \; ..., \; \lambda_n \geq 0$, что
        $\lambda_1 + ... + \lambda_n = 1$, выполняется неравенство
        $f(\lambda_1 x_1 + ... + \lambda_n x_n) \leq \lambda_1 f(x_1) + ... + \lambda_n f(x_n)$.
    \end{note}
\end{definition}

%% def 20
% Определение первообразной. 
% Формулировка основного свойства первообразных.
% !TeX root = ../colloquium-2.tex

% Определение первообразной. 
% Формулировка основного свойства первообразных.

\phantomsection\label{sec_Definition20}
\begin{definition}[
    Определение первообразной.
    Формулировка основного свойства первообразных.
]
    \subsection*{Определение первообразной.}
    \begin{note}[purple]
        Пусть функции $F$ и $f$ определены на интервале $(a, \; b)$. Если для любого числа
        $x \in (a, \; b)$ $F'(x) = f(x)$, то функция $F$ называется (\textit{точной})
        первообразной функции $f$ на интервале $(a, \; b)$. 
    \end{note}

    \subsection*{Формулировка основного свойства первообразных.}
    \begin{note}[purple]
        Пусть функции $F_1$ и $F_2$ являются первообразными функции $f$ на интервале $(a, \; b)$.
        Тогда существует такое число $C \in \mathbb{R}$, что $F_1(x) - F_2(x) = C \; 
        \forall x \in (a, \; b)$ 
    \end{note}
\end{definition}

%% def 21
% Определение неопределённого интеграла и таблица интегралов.
% !TeX root = ../colloquium-2.tex

% Определение неопределённого интеграла и таблица интегралов.

\phantomsection\label{sec_Definition21}
\begin{definition}[
    Определение неопределённого интеграла и таблица интегралов.
]
    \subsection*{Определение неопределённого интеграла.}
    \begin{note}[purple]
        Произвольная первообразная функциии $f$, определённой на некотором интервале, называется 
        неопределённым интегралом функции $f$ и обозначается $\int f(x) \, dx = F(x) + C$. \\
        "Произвольная первообразная функции $f$ равна какой-то известной первообразной плюс 
        некоторая постоянная".
    \end{note}

    \subsection*{Таблица интегралов.}
    \begin{note}[purple]
        \begin{enumerate}
            \item 
        \end{enumerate}
    \end{note}
\end{definition}

%% def 22
% Свойство линейности неопределённого интеграла. 
% Формулировка теоремы о формуле интегрирования по частям.
% Формулировка теоремы о замене переменной в неопределённом интервале.
% !TeX root = ../colloquium-2.tex

% Свойство линейности неопределённого интеграла. 
% Формулировка теоремы о формуле интегрирования по частям.
% Формулировка теоремы о замене переменной в неопределённом интервале.

\phantomsection\label{sec_Definition21}
\begin{definition}[
    Свойство линейности неопределённого интеграла. 
    Формулировка теоремы о формуле интегрирования по частям.
    Формулировка теоремы о замене переменной в неопределённом интервале.
]
    \subsection*{Свойство линейности неопределённого интеграла. }
    \begin{note}[purple]
        Если функции $F$ и $G$ дифференцируемы на интервале $(a, \; b)$ и $F' = f$,
        а $G' = g$, то в силу линейности производной при $\alpha, \beta \in \mathbb{R}$
        $$ (\alpha F + \beta G)' = \alpha f + \beta g $$
    \end{note}

    \subsection*{Формулировка теоремы о формуле интегрирования по частям.}
    \begin{note}[purple]
        Пусть функции $f$ и $g$ дифференцируемы на интервале $(a, \; b)$. Пусть также
        на этом интервале существует первообразная функции $f'g$. Тогда на интервале $(a, \; b)$
        существует первообразная функции $g'f$ и справедлива \textbf{формула интегрирования 
        по частям}. 
        $$ \int f'(x)g(x)dx = f(x)g(x) - \int f(x)g'(x)dx + C $$
    \end{note}

    \subsection*{Формулировка теоремы о замене переменной в неопределённом интервале.}
    \begin{note}[purple]
        Пусть функция $\varphi$ определена и дифференцируема на интервале $(\alpha, \; \beta)$
        и при любом $t \in (\alpha, \; \beta)$ значение $\varphi(t)$ принадлежит интервалу
        $(a, \; b)$, на котором определена функция $f$. Пусть также существует интеграл 
        $\int f(x)dx$. Тогда существует интеграл $\int f(\varphi(t)) \varphi'(t) dt$ и
        верна следующая формула замены переменной:
        $$ \int f(x)dx = \int f(\varphi(t)) \varphi'(t) dt $$
    \end{note}
\end{definition}

%% def 23
% Определение рациональной функции. 
% Четыре типа элементарных дробей. 
% Определение правильной дроби. 
% Метрика и метрическое пространство, норма и нормированное пространство.
% \input{definitions/definition-23.tex}

%% def 24
% Скалярное произведение и евклидово пространство.
% Неравенство Коши - Буняковского. 
% Связь нормы со скалярным произведением. 


%% def 25
% Евклидова метрика. 
% Шар, последовательность, предельная точка, открытое и замкнутое множества метрического пространства.


%% def 26
% Связь сходимости последовательностей и метрик и единственность предела. 
% Ограниченность. Арифметика пределов.
% Равносильность сходимости по норме и покоординатной сходимости.


%% def 27
% Предел функции по Коши.
% Предел функции по Гейне. 
% Равносильность определений предела по Коши и по Гейне.


%% def 28
% Единственность предела функции. 
% Арифметика пределов для фукнций. 
% Предел сложной функции. 
% Отделимость. 


%% def 29
% Непрерывность отображения по Гейне.
% Непрерывность отображения по Коши.
% Непрерывность сложной функции. 
% Локальные свойства непрерывных функций.


%% def 30
% Отображения координатных пространств. 
% Координатные функции. 
% Линейное отображение.
% Матрица линейного отображения.


%% def 31
% Единственность дифференциала. 
% Необходимое и достаточное условие дифференцируемости.